\documentclass[a4paper,10pt]{article}

\usepackage[margin=1in]{geometry} 	% Setea el margen manualmente, todos iguales.
\usepackage[spanish]{babel} 		% {Con estos dos anda
\usepackage[utf8]{inputenc} 		% todo lo que es tildes y ñ}
\usepackage{fancyhdr} 			%{Estos dos son para
\pagestyle{fancyplain} 			% el header copado}
\usepackage{color}			% Con esto puedo hacer la matufia de poner en color blanco un texto para engañar al formato
\usepackage{xcolor}
\usepackage{hyperref}
\usepackage{float}
\usepackage{framed}			%para crear cajas de texto
\usepackage{caratula}
\usepackage{verbatim}
\lhead{Paradigmas de Programación}% {Con esto se usa el header copado. También está \chead para
\rhead{Trabajo Práctico 1} 	% el centro y comandos para el pie de página, buscar fancyhdr}


%%%%%%%%%%%%%%%%%%%%%%%%%%%%%%%%%%%%%
%      COMANDOS ÚTILES USADOS       %
%%%%%%%%%%%%%%%%%%%%%%%%%%%%%%%%%%%%%

% \section{title} 		Te hace un título ``importante'' en negrita, numerado. También está \subsection{title} y \subsubsection{title}.
% \begin{itemize}		Te hace viñetas.
%	\item esto es un item	Cambiar itemize por enumerate te hace una numeración.
% \end{itemize}

% \textbf{text} 		Te hace el texto en negrita (bold).
% \underline{text}		Te subraya el texto.

% \textsuperscript{text}	Te hace ``superindices'' con texto. En teoría subscript debería funcionar, pero se puede usar guion bajo entre llaves
% 				y signos peso para hacerlo como alternativa. Sino buscar.

% \begin{tabular}{cols} 	Es para hacer tablas. Se pone una c por cada columna deseada dentro de cols (si es que se desea centrada, l para justificar a 
%	a & b & c		izquierda, r a la derecha). Si se separa por espacios la tabla no tendrá líneas divisorias. Si se separa por | en lugar de 
% \end{tabular}			espacios, aparecerá una línea. Con || dos, y así. Luego para los elementos de las filas se escriben y se separan con ampersand (&).
%				Finalmente, para las líneas horizontales, se usa \hline para una linea en toda la tabla y \cline{i - j} te hace la linea desde
%				la celda i hasta la j, arrancando en 1.
%				Si en la columna se pone p(width) podés escribir un párrafo en la celda. Para hacer un enter con \\ no funciona porque te hace un
%				enter en la fila. Para eso se usa el comando \newline.
  
% \textcolor{color predefinido en palabras}{text}

%%%%%%%%%%%%%%%%%%%%%%%%%%%%%%%%%%%%%
%    FIN COMANDOS ÚTILES USADOS     %
%%%%%%%%%%%%%%%%%%%%%%%%%%%%%%%%%%%%%

\begin{document}

%%%%%%%%%%%%%%%%%%%%%%%%%%%%%%
%    Carátula										      %
%%%%%%%%%%%%%%%%%%%%%%%%%%%%%%



% **************************************************************************
%
%  Package 'caratula', version 0.3 (para componer caratulas de TPs del DC).
%
%  En caso de dudas, problemas o sugerencias sobre este package escribir a
%  Brian J. Cardiff (bcardif arroba gmail.com).
%  Nico Rosner (nrosner arroba dc.uba.ar).
%
% **************************************************************************

% ----- Informacion sobre el package para el sistema -----------------------

\NeedsTeXFormat{LaTeX2e}
\ProvidesPackage{caratula}[2005/08/09 v0.3 Para componer caratulas de TPs del DC]
\usepackage[pdftex]{graphicx}

% ----- Imprimir un mensajito al procesar un .tex que use este package -----

\typeout{Cargando package 'caratula' v0.3 (2005/08/09)}

% ----- Algunas variables --------------------------------------------------

\let\Materia\relax
\let\Submateria\relax
\let\Titulo\relax
\let\Subtitulo\relax
\let\Grupo\relax
\let\Fecha\relax
\let\Logoimagefile\relax
\let\Resumen\relax
\let\Tags\relax

% ----- Comandos para que el usuario defina las variables ------------------

\def\materia#1{\def\Materia{#1}}
\def\submateria#1{\def\Submateria{#1}}
\def\titulo#1{\def\Titulo{#1}}
\def\subtitulo#1{\def\Subtitulo{#1}}
\def\grupo#1{\def\Grupo{#1}}
\def\fecha#1{\def\Fecha{#1}}
\def\logoimagefile#1{\def\Logoimagefile{#1}}
\def\resumen#1{\def\Resumen{#1}}
\def\tags#1{\def\Tags{#1}}

% ----- Token list para los integrantes ------------------------------------

\newtoks\intlist\intlist={}

% ----- Comando para que el usuario agregue integrantes --------------------

\def\integrante#1#2#3{\intlist=\expandafter{\the\intlist
    \rule{0pt}{1.2em}#1&#2&\tt #3\\[0.2em]}}

% ----- Macro para generar la tabla de integrantes -------------------------

\def\tablaints{%
    \begin{tabular}[t]{| l @{\hspace{4ex}} c @{\hspace{4ex}} l|}
        \hline
        \multicolumn{1}{|c}{\rule{0pt}{1.2em} Integrante} & LU &  \multicolumn{1}{c|}{Correo electr\'onico} \\[0.2em]
        \hline \hline
        \the\intlist
        \hline
    \end{tabular}}

% ----- Codigo para manejo de errores --------------------------------------

\def\se{\let\ifsetuperror\iftrue}
\def\ifsetuperror{%
    \let\ifsetuperror\iffalse
    \ifx\Materia\relax\se\errhelp={Te olvidaste de proveer una \materia{}.}\fi
    \ifx\Titulo\relax\se\errhelp={Te olvidaste de proveer un \titulo{}.}\fi
    \edef\mlist{\the\intlist}\ifx\mlist\empty\se%
    \errhelp={Tenes que proveer al menos un \integrante{nombre}{lu}{email}.}\fi
    \expandafter\ifsetuperror}


% ----- \maketitletxt correspondiente a la versión v0.2 (texto) ---------

\def\maketitletxt{%
    \ifsetuperror\errmessage{Faltan datos de la caratula! Ingresar 'h' para mas informacion.}\fi
    \thispagestyle{empty}
    \begin{center}
    \vspace*{\stretch{2}}
    {\LARGE\textbf{\Materia}}\\[1em]
    \ifx\Submateria\relax\else{\Large \Submateria}\\[0.5em]\fi
    \par\vspace{\stretch{1}}
    {\large Departamento de Computaci\'on}\\[0.5em]
    {\large Facultad de Ciencias Exactas y Naturales}\\[0.4em]
    {\large Universidad de Buenos Aires}
    \par\vspace{\stretch{3}}
    {\Large \textbf{\Titulo}}\\[0.8em]
    {\Large \Subtitulo}
    \par\vspace{\stretch{3}}
    \ifx\Grupo\relax\else\textbf{\Grupo}\par\bigskip\fi
    \tablaints
    \end{center}
    \vspace*{\stretch{3}}
    \newpage}

% ----- \maketitle correspondiente a la versión v0.3 (gráfica) -------------

\def\maketitlegraf{%
    \ifsetuperror\errmessage{Faltan datos de la caratula! Ingresar 'h' para mas informacion.}\fi
%
    \thispagestyle{empty}

    \ifx\Logoimagefile\relax\else\includegraphics{\Logoimagefile}\fi \hfill \includegraphics{logo_dc.jpg}

    \vspace*{.07 \textheight}

    \noindent \textbf{\huge \Titulo}  \medskip \\
    \ifx\Subtitulo\relax\else\noindent\textbf{\large \Subtitulo} \\ \fi%
    \noindent \rule{\textwidth}{1 pt}

    {\noindent\large\Fecha \hspace*\fill \Materia} \\
    \ifx\Submateria\relax\else{\noindent \hspace*\fill \Submateria}\fi%

    \medskip%
    \begin{center}
        \ifx\Grupo\relax\else\textbf{\Grupo}\par\bigskip\fi
        \tablaints
    \end{center}%
    \bigskip
    %~ \ifx\Resumen\relax\else\textbf{\textsf{Resumen}} \\
    %~ \Resumen\bigskip\fi \\
    %~ \ifx\Tags\relax\else\textbf{\textsf{Palabras Clave}} \\
    %~ \Tags\par\bigskip\fi
    \vfill%
%
    \begin{minipage}[t]{\textwidth}
        \begin{minipage}[t]{.50 \textwidth}
            \includegraphics[width=0.9\textwidth]{logo_uba.jpg}
        \end{minipage}%%
        \begin{minipage}[b]{.50 \textwidth}
            \textbf{\textsf{Facultad de Ciencias Exactas y Naturales}} \\
            \textsf{Universidad de Buenos Aires} \\
            {\scriptsize %
            Ciudad Universitaria - (Pabell\'on I/Planta Baja) \\
                Intendente G\"uiraldes 2160 - C1428EGA \\
            Ciudad Aut\'onoma de Buenos Aires - Rep. Argentina \\
                Tel/Fax: (54 11) 4576-3359 \\
            http://www.fcen.uba.ar \\
            }
        \end{minipage}
    \end{minipage}%
%
    \newpage}

% ----- Reemplazamos el comando \maketitle de LaTeX con el nuestro ---------

\def\maketitle{\maketitlegraf}
\newpage

%%%%%%%%%%%%%%%%%%%%%%%%%%%%%%%%%%%%
%    Tabla de contenidos/Índice    						%
%%%%%%%%%%%%%%%%%%%%%%%%%%%%%%%%%%%%

\setcounter{page}{1}

%-- Indice --
\newpage{\pagestyle{empty}\tableofcontents\cleardoublepage}

\section{Tipos}

\begin{verbatim}

module Tipos where

import Grafo

type Mundo = Integer
type Prop = String
data Modelo = K (Grafo Mundo) (Prop -> [Mundo])

data Exp = Var Prop | Not Exp | Or Exp Exp | And Exp Exp | D Exp | B Exp deriving (Show, Eq)

\end{verbatim}

\newpage

\section{Grafo}

\begin{verbatim}

module Grafo 
    (Grafo, vacio, nodos, vecinos, agNodo, sacarNodo, agEje, lineal, union, clausura) where

import qualified Data.List as List (delete, union, nub)

-- Invariante: ninguna de las listas (nodos y vecinos de nodos) tienen repetidos
data Grafo a = G [a] (a -> [a])

instance (Show a) => Show (Grafo a) where
  show (G n e) = 
    "[\n" ++ concat (map (\x -> " " ++ show x ++ " -> " ++ show (e x) ++ "\n") n) ++ "]"


-- Ejercicio 1
-- Crea un grafo sin nodos.

vacio :: Grafo a
vacio = G [] (\_ -> [])

-- Ejercicio 2
-- Devuelve los nodos de un grafo.

nodos :: Grafo a -> [a]
nodos (G ns fv) = ns

-- Ejercicio 3
-- Devuelve los vecinos de un nodo en un grafo dado.

vecinos :: Grafo a -> a -> [a]
vecinos (G ns fv) = fv

-- Ejercicio 4
-- Devuelve un grafo con un nodo agregado, sin vecinos.
-- Si el nodo ya está presente en el grafo, no tendrá vecinos en nuevo grafo.

agNodo :: Eq a => a -> Grafo a -> Grafo a
agNodo n (G ns fv) = G ns' fv'
        where ns' = List.union [n] ns
              fv' e | e == n    = []
                    | otherwise = fv e


-- Ejercicio 5
-- Devuelve un grafo sin un nodo dado.

sacarNodo :: Eq a => a -> Grafo a -> Grafo a
sacarNodo n (G ns fv) = G ns' fv'
        where ns' = List.delete n ns
              fv' e | e == n    = []
                    | otherwise = List.delete n $ fv e






-- Ejercicio 6
-- Devuelve un grafo con un eje agregado.

agEje :: Eq a =>  (a,a) -> Grafo a -> Grafo a
agEje (n1, n2) (G ns fv) = G ns fv' 
        where fv' e | e == n1   = n2 : fv e
                    | otherwise = fv e


-- Ejercicio 7
-- Devuelve un grafo donde los nodos son todos los de la lista pasada por
-- argumento, y cada nodo tiene como unico vecino al elemento que lo sigue en
-- dicha lista.

lineal :: Eq a => [a] -> Grafo a
lineal ns = G ns fv 
        where fv e = take 1 $ drop 1 $ dropWhile (/= e) ns


-- Ejercicio 8
-- Devuelve la unión de dos grafos.

union :: Eq a => Grafo a -> Grafo a -> Grafo a
union (G ns1 fv1) (G ns2 fv2) = G ns' fv'
        where ns' = List.union ns1 ns2
              fv' e = List.union (fv1 e) (fv2 e)

-- Ejercicio 9
-- Devuelve la clausura reflexo transitiva de un grafo.

clausura :: (Eq a) => Grafo a -> Grafo a
clausura (G ns fv) = G ns fv'
        where fv' e = puntofijo extenderConVecinos [e]
              extenderConVecinos xs = foldl List.union xs (map fv xs)

-- Punto fijo
-- Devuelve f(f(...f(x)...)) = f^n(x), tal que f(f^n(x)) = f^n(x), es decir que
-- es un punto fijo de la función.

puntofijo :: (Eq a) => (a -> a) -> (a -> a)
puntofijo f x = head $ dropWhile (\e -> (f e) /= e) $ iterate f x

\end{verbatim}

\newpage

\section{Lomoba}


\begin{verbatim}

module Lomoba where
import Grafo
import Tipos
import qualified Data.List as List (union, intersect)


-- Ejercicio 10
-- Realiza recursión estructural sobre expresiónes.

foldExp :: 
  (Prop -> b) -> (b -> b) -> (b -> b -> b) -> (b -> b -> b) -> (b -> b) -> (b -> b) -> Exp -> b
foldExp fVar fNot fOr fAnd fD fB ei =
        let frec = foldExp fVar fNot fOr fAnd fD fB
        in case ei of
                Var p -> fVar p
                Not e -> fNot (frec e)
                Or e1 e2 -> fOr (frec e1) (frec e2)
                And e3 e4 -> fAnd (frec e3) (frec e4)
                D e5 -> fD (frec e5)
                B e6 -> fB (frec e6)


-- Ejercicio 11
-- Devuelve un la visibilidad de una fórmula, es decir cuánto del grafo se
-- utiliza para evaluarla.

visibilidad :: Exp -> Integer
visibilidad = foldExp fVar fNot fOr fAnd fD fB
        where   fVar = const 0
                fNot = id
                fOr = max
                fAnd = max
                fD = (1+)
                fB = (1+)


-- Ejercicio 12
-- Lista las variables que aparecen en la fórmula.

extraer :: Exp -> [Prop]
extraer = foldExp fVar fNot fOr fAnd fD fB
        where   fVar p = [p]
                fNot = id
                fOr = List.union
                fAnd = List.union
                fD = id
                fB = id

-- Ejercicio 13
-- Dado un modelo, decide si una fórmula es verdadera en un mundo.

eval :: Modelo -> Mundo -> Exp -> Bool
eval m w e = eval' m e w


eval' :: Modelo -> Exp -> Mundo -> Bool
eval' m@(K g v) = foldExp fVar fNot fOr fAnd fD fB
        where        fVar x w = elem w (v x)
                fNot rec w = not (rec w)
                fOr recI recD w = (recI w) || (recD w)
                fAnd recI recD w = (recI w) && (recD w)
                fD rec w = any rec (vecinos g w)
                fB rec w = all rec (vecinos g w)


-- Ejercicio 14
-- Devuelve todos los mundos de un modelo para los que vale una expresión.

valeEn :: Exp -> Modelo -> [Mundo]
valeEn e m@(K g v) = filter (eval' m e) (nodos g)


-- Ejercicio 15
-- Devuelve un modelo en el que para todos los mundos vale la expresión dada.

quitar :: Exp -> Modelo -> Modelo
quitar e m@(K g v) = K g' v'
        where   g' = foldl (flip sacarNodo) g (noValeEn e m)
                v' p = List.intersect (nodos g') (v p)

noValeEn :: Exp -> Modelo ->[Mundo]
noValeEn e m@(K g v) = filter (not . eval' m e) (nodos g)

-- Ejercicio 16
-- Dado un modelo, indica si para todos sus mundos vale una expresión.

cierto :: Modelo -> Exp -> Bool
cierto m@(K g v) e = all (eval' m e) (nodos g)

\end{verbatim}

%\newpage

%\section{Tests}

\begin{verbatim}
import Grafo
import Tipos
import Lomoba
import Parser
import Test.HUnit

-- evaluar t para correr todos los tests
t = runTestTT allTests

allTests = test [
	"parser" ~: testsParser,
	"grafo" ~: testsGrafo,
	"lomoba" ~: testsLomoba
	]

testsParser = test [
	(Var "p") 						~=? (parse "p"),
	(And (Var "p") (Var "q")) 		~=? (parse "p && q"),
	(Or (Var "p") (Var "q")) 		~=? (parse "p || q"),
	(Or (Not (Var "p")) (Var "q"))	~=? (parse "!p || q"),
	(And (D (Var "p")) (Var "q")) 	~=? (parse "<>p && q"),
	(And (B (Var "p")) (Var "q")) 	~=? (parse "[]p && q"),
	(D (And (Var "p") (Var "q"))) 	~=? (parse "<>(p && q)"),
	(B (And (Var "p") (Var "q"))) 	~=? (parse "[](p && q)")]

testsGrafo = test [
	--nodos	y agNodo
	[1] ~~? (nodos (agNodo 1 vacio)),
	[1,2] ~~? (nodos (agNodo 2 (agNodo 1 vacio))),

	--vecinos y egEje
	[1] ~~? (vecinos (agEje (2,1) (agNodo 2 (agNodo 1 vacio))) 2 ),
	[] ~~? (vecinos (agEje (2,1) (agNodo 2 (agNodo 1 vacio))) 1 ),
	[1,2,3] ~~? (nodos((agEje(3,2) (agEje (1,3) (agEje (1,2)(agNodo 3 (agNodo 2 (agNodo 1 vacio)))))))),
	[2,3] ~~? (vecinos (agEje(3,2) (agEje (1,3) (agEje (1,2)(agNodo 3 (agNodo 2 (agNodo 1 vacio)))))) 1),
	[2] ~~? (vecinos (agEje(3,2) (agEje (1,3) (agEje (1,2)(agNodo 3 (agNodo 2 (agNodo 1 vacio)))))) 3),
	[] ~~? (vecinos (agEje(3,2) (agEje (1,3) (agEje (1,2)(agNodo 3 (agNodo 2 (agNodo 1 vacio)))))) 2),	

	--sacarNodo
	[2] ~~? (nodos(sacarNodo 1 ((agNodo 2 (agNodo 1 vacio))))),	
	[] ~~? (vecinos (sacarNodo 2 (agEje(3,2) (agEje (1,3) (agEje (1,2)(agNodo 3 (agNodo 2 (agNodo 1 vacio))))))) 3 ),
	[3] ~~? (vecinos (sacarNodo 2 (agEje(3,2) (agEje (1,3) (agEje (1,2)(agNodo 3 (agNodo 2 (agNodo 1 vacio))))))) 1 ),
	[1,3] ~~? (nodos (sacarNodo 2 (agEje(3,2) (agEje (1,3) (agEje (1,2)(agNodo 3 (agNodo 2 (agNodo 1 vacio)))))))),	

	--lineal	
	[1,2,3,4] ~~? (nodos (lineal [1,2,3,4])), 	
	[2] ~~?	(vecinos (lineal [1,2,3,4]) 1),
	[3] ~~?	(vecinos (lineal [1,2,3,4]) 2),
	[4] ~~?	(vecinos (lineal [1,2,3,4]) 3),
	[] ~~?	(vecinos (lineal [1,2,3,4]) 4),

	--union
	[1,2,3] ~~? (nodos (union (agEje (1,2) (agNodo 2 (agNodo 1 vacio))) (agEje (3,1) (agEje (1,3) (agNodo 3 (agNodo 1 vacio)))))),
	[2,3] ~~? (vecinos (union (agEje (1,2) (agNodo 2 (agNodo 1 vacio))) (agEje (3,1) (agEje (1,3) (agNodo 3 (agNodo 1 vacio))))) 1),
	[] ~~? (vecinos (union (agEje (1,2) (agNodo 2 (agNodo 1 vacio))) (agEje (3,1) (agEje (1,3) (agNodo 3 (agNodo 1 vacio))))) 2),
	[1] ~~? (vecinos (union (agEje (1,2) (agNodo 2 (agNodo 1 vacio))) (agEje (3,1) (agEje (1,3) (agNodo 3 (agNodo 1 vacio))))) 3),
	[1,2,3,4] ~~? (nodos (union (vacio) (lineal [1,2,3,4]))),	
	
	--clasura	
	[1,2,3,4] ~~? (nodos (clausura (lineal [1,2,3,4]))),
	[1,2,3,4] ~~? (vecinos (clausura (lineal [1,2,3,4])) 1),
	[2,3,4] ~~? (vecinos (clausura (lineal [1,2,3,4])) 2),
	[3,4] ~~? (vecinos (clausura (lineal [1,2,3,4])) 3),
	[4] ~~?(vecinos (clausura (lineal [1,2,3,4])) 4)
	]

testsLomoba = test [
	-- foldExp
        exp1 ~=? (foldExp Var Not Or And D B exp1),

	-- visibilidad
	0 ~=? (visibilidad (parse "p")),
	1 ~=? (visibilidad (parse "<>p")),
	1 ~=? (visibilidad (parse "!<>p")),
	2 ~=? (visibilidad (parse "<>!<>p")),
	2 ~=? (visibilidad (parse "<><>p || <><>q")),
	3 ~=? (visibilidad (parse "<>(<>p || <><>q)")),
	3 ~=? (visibilidad (parse "[](<>p && <>[]q)")),

	-- extraer
	["p"] ~~? (extraer (parse "p")),
	["p"] ~~? (extraer (parse "<>p")),
	["p"] ~~? (extraer (parse "[]p")),
	["p", "q"] ~~? (extraer (parse "p||q")),	
	["p", "q"] ~~? (extraer (parse "p&&q")),
	["p"] ~~? (extraer (parse "<><>p || <><>p")),
	["p", "q", "r"] ~~? (extraer (parse "(p||q)&&[]<>r")),	
	["p", "q", "r"] ~~? (extraer (parse "<>((p||q)&&[]<>r)")),

	-- eval
	True ~=? eval modelo1 1 (parse "p"),
	False ~=? eval modelo1 1 (parse "q"),

	True ~=? eval modelo1 4 (parse "r && q"),
	False ~=? eval modelo1 4 (parse "p && q"),

	True ~=? eval modelo1 1 (parse "p || q"),
	False ~=? eval modelo1 2 (parse "p || q"),

	True ~=? eval modelo1 2 (parse "!p || q"),
	False ~=? eval modelo1 1 (parse "!p || q"),
	True ~=? eval modelo1 3 (parse "!p || q"),

	False ~=? eval modelo1 1 (parse "<>p && p"),
	True ~=? eval modelo1 1 (parse "<>q && p"),

	True ~=? eval modelo1 1 (parse "[]r"),
	False ~=? eval modelo1 1 (parse "[]q"),

	False ~=? eval modelo1 1 (parse "<>(p && q)"),
	True ~=? eval modelo1 1 (parse "<>(r && q)"),

	True ~=? eval modelo1 2 (parse "[](!p && q)"),

	-- valeEn
	[1] ~~? valeEn (parse "p") modelo1,
	[2,3,4] ~~? valeEn (parse "!p") modelo1,
	[3,4] ~~? valeEn (parse "q") modelo1,
	[1,2] ~~? valeEn (parse "<>q") modelo1,
	[2,3,4] ~~? valeEn (parse "[]q") modelo1,
	[1,3,4] ~~? valeEn (parse "[]r") modelo1,
	[4] ~~? valeEn (parse "r && q") modelo1,
	[1,2,3,4] ~~? valeEn (parse "r || (q || p)") modelo1,
	[1,3] ~~? valeEn (parse "!r && (q || p)") modelo1,
	[] ~~? valeEn (parse "s") modelo1,
	
	
	--quitar
	[1] ~~? let (K g _) = quitar (parse "p") modelo1 in nodos g,
	[4,2] ~~? let (K g _) = quitar (parse "r") modelo1 in nodos g,
	[4] ~~? let (K g _) = quitar (parse "r&&q") modelo1 in nodos g,
	[1,3,4] ~~? let (K g _) = quitar (parse "[]r") modelo1 in nodos g,
	[1] ~~? let (K g _) = quitar (parse "<>r") modelo1 in nodos g,

	--cierto
	True ~=? cierto (K vacio (const [])) (parse "p"),
	True ~=? cierto (K vacio (const [])) (parse "[]p"),
	False ~=? cierto modelo1 (parse "p"),
	True ~=? cierto modelo1 (parse "[][]q")
	]

modelo1 = K (union (lineal [1,2,3]) (lineal [1,4]) ) v1
	where	v1 "p" = [1] 
		v1 "q" = [3,4]
		v1 "r" = [2,4]
		v1 _ = []

exp1 = parse "<>[]p || q && r"

---------------
--  helpers  --
---------------

-- idem ~=? pero sin importar el orden
(~~?) :: (Ord a, Eq a, Show a) => [a] -> [a] -> Test
expected ~~? actual = (sort expected) ~=? (sort actual)
	where
		sort = foldl (\r e -> push r e) []
		push r e = (filter (e<=) r) ++ [e] ++ (filter (e>) r)
\end{verbatim}

%\newpage





%%%%%%%%%%%%%%%%%%%%%%%%%%%%%%%%%%%%%%%%
%        Esto es de mi template        %
%    Dejar para futuras referencias    %
%%%%%%%%%%%%%%%%%%%%%%%%%%%%%%%%%%%%%%%%

% \section{Sección, título grande}
% \\
% \textcolor{white}{Sarasa engañadora de formato, jua jua}\\
% Esto es texto. Con doble barra invertida es el enter.
% \\
% \\
% \textit{De acá en adelante, todo lo que se explica como "para hacer x se usa y", implica que antes de y va una barra invertida}
% \\
% \\
% Con \textbf{textbf\{\}} se escribe en negrita.
% \\
% \\
% \section{Bullet \& Numbering}
% Para hacer items se usa \textbf{begin\{itemize\}} y \textbf{end\{itemize\}}, por ejemplo:
% \begin{itemize}
% 	\item Esto es un ítem.
% 	\item Esto es otro.
% \end{itemize}
% Con \textbf{item} se hace cada uno de los ítems. Si cambiás \textbf{itemize} por \textbf{enumerate} te lo hace enumerado. Por ejemplo:
% \begin{enumerate}
% 	\item Acá está el primero.
% 	\item Acá el segundo.
% \end{enumerate}
% \\
% Y así.
% \\
% \\
% \section{Tablas}
% También tenemos las tablas, que son un poco más rebuscadas. Se usa \textbf{begin\{tabular\}\{cols\}} y en cols ponemos c si queremos una columna centrada, l y r
% para otro tipo de justificación. Si las c las separás con espacios, se hacen columnas sin división. Si ponés un pipe es con una línea divisoria, dos pipes con
% dos líneas, y así. Se termina con \textbf{end\{tabular\}}. Para separar entre elementos de fila/columna se usa un ampersand (\&, y es necesario para separar
% elementos entre filas y columnas, no solo entre filas) y para cambiar de fila \textbf{SIN} linea divisoria, un \textbf{newline} ya que el doble barra invertida
% te hace un enter dentro de la celda. Con línea divisoria es reemplazando \textbf{newline} por \textbf{hline}.
% \\
% \\
% Más datos en el principio de este tex. Un ejemplo:
% \\
% \\
% \begin{tabular}{| c| c|}\hline
%     Celda 1 & Celda 2 &\hline
%     Celda 3 & Celda 4 &\hline
% \end{tabular}
% \\
% \\
% \section{Verbatim}
% \begin{verbatim}
% Esto es verbatim. Es un entorno que no le da ni 5 de pelota al formato de latex.
% Por eso mismo hay que tener cuidado con no irse de la hoja o similar.
% Sirve por ejemplo, para pseudocódigo:
% 
% if (se cumple sarasa){
%     ejecuto cosito1;
%     ejecuto cosito2:
% }else{
%     ejecuto cosito3:
% }
% \end{verbatim}

\end{document}
